\section{Vehicles}

\subsection{Basilisk}

Basilisks can be relaid, reloaded, and fired in a relatively short time, and the design of the Earthshaker Cannon allows it to be elevated high enough to fire its shells in an arc towards concealed targets. They can thus be deployed out of the enemy's reach, sheltered from any return fire. Its lack of armour, coupled with an open platform that exposes the crew, makes it too vulnerable to enemy fire during an assault.

The Basilisk has the \textbf{Artillery} and \textbf{Slave (Nest)} abilities.

\begin{unitstatstable}
15 cm & 4+ & +0 & 8 & Earthshaker Cannon & 150 cm & 2 & 4+ & $-1$
\end{unitstatstable}

\unitimage{basilisk.png}

\unitdivider

\subsection{Bombard}

Bombards are among the oldest and most famous artillery pieces of the Imperium, and are among the largest used routinely by the Imperial Guard's heavy artillery. They are seen only rarely, however, except during campaigns in which their siege mortars---of unprecedented calibre---are required to atomise enemy positions.

The Bombard uses a \textbf{Template (7.5 cm)} and has the \textbf{Artillery}, \textbf{Battery}, \textbf{Slave (Nest)}, and \textbf{Damages Buildings (AP $-3$/2)} abilities.

\begin{unitstatstable}
15 cm & 4+ & +0 & 8 & Colossus Mortar & 45--150 cm & Template & 3+ & $-2$
\end{unitstatstable}

\unitimage{bombard.png}

\unitdivider

\subsection{Goliath Truck}

Goliaths are robust transports originally designed to carry Imperial overseers through underground complexes and mining tunnels. The vehicle's dense construction can withstand the most hostile subterranean environments, and it is covered in layers of chemically treated permacite to offer a measure of protection against every industrial accident the Imperium has already faced. An ascending cult will use every type of vehicle, from lunar quads to stretched civilian cars, not forgetting mobile industrial macro-platforms, but the Goliath remains the most prized.

It has the \textbf{Transport (2)}, \textbf{Attached Transport}, and \textbf{Slave (Hunt)} abilities.

\begin{unitstatstable}
25 cm & 5+ & +1 & Attached & Autocannon & 45 cm & 1 & 4+ & 0
\end{unitstatstable}

\unitimage{goliath-truck.png}

\unitdivider

\subsection{Hellhound}

Many Flamer weapons tend to suffer from a lack of ammunition capacity, requiring a considerable quantity of fuel to maintain a regular rate of fire. The Hellhound circumvents the problem by sacrificing the Chimera's transport capacity in favour of more powerful engines and prometheum tanks, feeding the monstrous turret Flamer---the aptly named Inferno Cannon---which projects a spontaneously combusting fluid. Generally deployed in dense terrain, the Hellhound's principal role is as a terror weapon during intense engagements, such as street fighting, or against enemies concealed---and easily flushed out---in bunkers and other earthworks.

Both of the Hellhound's weapons have the \textbf{Reduces Cover ($-3$)} ability, and its Inferno Cannon has the \textbf{Turret} ability. In addition, the Hellhound has the \textbf{Slave (Hunt)} ability.

\begin{unitstatstable}
\SetCell[r=2]{c,m} 20 cm &
\SetCell[r=2]{c,m} 3+ &
\SetCell[r=2]{c,m} +0 &
\SetCell[r=2]{c,m} 8 &
Inferno Cannon &
20 cm &
2 &
4+ &
0
\\
& & & &
Flamer &
20 cm &
1 &
4+ &
0
\end{unitstatstable}

\unitimage{hellhound.png}

\unitdivider

\subsection{Griffon}

The Griffon is equipped with a heavy mortar to which its experienced crews can give a high rate of fire. This vehicle is designed to serve as mobile infantry support at short and medium range. It can maintain a high rate of fire and has effective firepower against infantry and light vehicles, allowing heavier weapons to concentrate on tougher targets. Well used, the Griffon is perfectly suited both to defensive fire and to close support during offensive operations, and it is a tactically versatile weapon when it employs different types of ammunition.

The Griffon has the \textbf{Artillery}, \textbf{Slave (Nest)}, \textbf{Battery}, and \textbf{Reduces Cover ($-1$)} abilities, and uses a \textbf{Template (12 cm)}.

\begin{unitstatstable}
15 cm & 4+ & +0 & 8 & Griffon Heavy Mortar & 150 cm & Template & 4+ & 0
\end{unitstatstable}

\unitimage{griffon.png}

\unitdivider

\subsection{Goliath Rockgrinder}

The most widespread variant of the Goliath is the Rockgrinder. The enormous crushing blade at the front of the Rockgrinder possesses series of grinders capable of milling veins of ore and driving tunnels through bedrock. Their use is far more sinister in time of war. When the order is given to abandon all caution in favour of a general attack, the Rockgrinder's driver will race towards the front line, its gunner projecting tongues of blazing prometheum to keep the enemy's head down until impact.

It has the \textbf{Roller (5+ / AP $-1$)}, \textbf{Transport (2)}, \textbf{Attached Transport}, and \textbf{Slave (Devastation)} abilities.

\begin{unitstatstable}
25 cm & 4+ & +3 & Attached & Multi-melta & 20 cm & 1 & 3+ & $-2$
\end{unitstatstable}

\unitimage{goliath-rockgrinder.png}

\unitdivider

\subsection{Hellhound Devil Dog}

On the Devil Dog, the Hellhound's principal armament gives way to a squat fusion cannon. The tank's name comes from the high-pitched hiss this weapon produces, which allows the Devil Dog to engage targets far more heavily armoured than itself. Indeed, Devil Dog crews often see themselves as big-game hunters, practised in the pursuit of tanks as well as exceptionally resilient infantry. An enemy engaged in a prolonged campaign in dense terrain against the Astra Militarum learns to dread the compact silhouette of the Devil Dog, and those who do not feel this salutary fear find their vehicles reduced to slag by this fast and aggressive tank hunter.

Its Melta Cannon has the \textbf{Turret} ability, while its Flamer has the \textbf{Reduces Cover ($-3$)} ability, and the Devil Dog has the \textbf{Slave (Hunt)} ability.

\begin{unitstatstable}
\SetCell[r=2]{c,m} 20 cm &
\SetCell[r=2]{c,m} 3+ &
\SetCell[r=2]{c,m} +0 &
\SetCell[r=2]{c,m} 8 &
Melta Cannon &
20 cm &
1 &
3+ &
$-1$
\\
& & & &
Flamer &
20 cm &
1 &
4+ &
0
\end{unitstatstable}

\unitimage{hellhound-devil-dog.png}

\unitdivider

\subsection{Hydra}

Whether facing lightning-fast xenos fighters or daemonic monstrosities with membranous wings, the Hydra anti-aircraft tank possesses the ideal armament to empty the skies of airborne enemies. The Hydra's predictive logic spirits detect and lock onto targets with the tenacity of a predator. The autoloaders engage as the turret pivots with a roar and its quad Autocannons fill the air with tracer shells. Few aerial targets, supernatural or otherwise, survive once they find themselves in a Hydra's sights.

The Hydra has the \textbf{AA} and \textbf{Slave (Nest)} abilities.

\begin{unitstatstable}
15 cm & 4+ & +0 & 8 & Hydra Autocannon & 75 cm & 4 & 5+ & $-1$
\end{unitstatstable}

\unitimage{hydra.png}

\unitdivider

\subsection{Leman Russ}

The simplicity of this tank makes it easy to manufacture and maintain, so that it is produced by the million upon Forge Worlds and in the hive factories of the entire galaxy, and Imperial armies can field several dozen companies of Leman Russes in order to overwhelm the adversary, thus providing the regiments of the Imperial Guard and the Planetary Defence Forces with a powerful armoured support indispensable for survival on the brutal and implacable battlefields of the 41st Millennium. Only the most powerful shots succeed in penetrating its thick armour, allowing this vehicle to defend strategic positions, while its varied armament allows it to punch through enemy lines with ease when it goes on the offensive. Its simple and robust design has allowed it to stand the test of centuries, while its versatility and ease of construction make it one of the tanks most in demand among Guard officers.

Its Battle Cannon has the \textbf{Turret} ability, and the Leman Russ has the \textbf{Slave (Hunt)} ability.

\begin{unitstatstable}
\SetCell[r=3]{c,m} 15 cm &
\SetCell[r=3]{c,m} 2+ &
\SetCell[r=3]{c,m} +3 &
\SetCell[r=3]{c,m} 8 &
Battle Cannon &
75 cm &
1 &
4+ &
$-2$
\\
& & & &
Lascannon &
60 cm &
1 &
5+ &
$-1$
\\
& & & &
Heavy Bolters &
45 cm &
2 &
5+ &
$-1$
\end{unitstatstable}

\unitimage{leman-russ.png}

\unitdivider

\subsection{Leman Russ Annihilator}

This Leman Russ variant exchanges its battle cannon for Twin Lascannons, and this change of firepower makes the Annihilator a very effective tank hunter. The Annihilator is favoured in war zones where the rate of tank attrition is high and resupply poses a problem. As a result, the powerful foundries of Mars and the principal Forge Worlds of Ryza and Accatran have begun to produce the Leman Russ Annihilator in increasing numbers to meet demand following battlefield losses, and a number of other Forge Worlds across the galaxy are pressing for unrestricted access to the tank's plans.

Its Twin Lascannons have the \textbf{Turret} ability, and the Leman Russ Annihilator has the \textbf{Slave (Hunt)} ability.

\begin{unitstatstable}
\SetCell[r=2]{c,m} 15 cm &
\SetCell[r=2]{c,m} 2+ &
\SetCell[r=2]{c,m} +3 &
\SetCell[r=2]{c,m} 8 &
Twin Lascannons &
60 cm &
1 &
4+ &
$-1$
\\
& & & &
Lascannon &
60 cm &
1 &
5+ &
$-1$
\end{unitstatstable}

\unitimage{leman-russ-annihilator.png}

\unitdivider

\subsection{Leman Russ Demolisher}

The Leman Russ Demolisher is a heavily armoured mobile siege tank, used almost exclusively when Command seeks a breakthrough. This tank was built for one role: to break the lines. Thanks to its reinforced frontal armour, the Demolisher is able to advance under heavy fire in order to bring its cannon within range of an enemy fortress or other fortification. Its Demolisher Cannon possesses staggering destructive power, allowing it to dominate every short-range engagement, and can explode all but the most heavily reinforced structures, with the result that a breach is created through which armour and other infantry can fall upon those who have somehow managed to withstand the Demolisher's devastating barrage.

Its Demolisher Cannon has the \textbf{Turret} and \textbf{Reduces Cover ($-3$)} abilities, while its Flamer has the \textbf{Reduces Cover ($-3$)} ability, and the Leman Russ Demolisher has the \textbf{Slave (Hunt)} ability.

\begin{unitstatstable}
\SetCell[r=2]{c,m} 15 cm &
\SetCell[r=2]{c,m} 2+ &
\SetCell[r=2]{c,m} +3 &
\SetCell[r=2]{c,m} 8 &
Demolisher Cannon &
45 cm &
1 &
4+ &
$-3$
\\
& & & &
Flamer &
20 cm &
1 &
4+ &
0
\end{unitstatstable}

\unitimage{leman-russ-demolisher.png}

\unitdivider

\subsection{Leman Russ Exterminator}

The Leman Russ Exterminator is a widespread variant of the standard pattern, replacing its battle cannon with Twin Autocannons. Generally equipped with multiple Heavy Bolters mounted in turrets and on the hull for additional firepower, the Leman Russ Exterminator can unleash a true uninterrupted hail of projectiles that easily shred light armour and infantry within range of its crew's guns. In particular, by forgoing its anti-tank effectiveness for increased anti-infantry firepower, the Exterminator comes into its own when deployed against hordes of lightly equipped enemy troops such as Ork mobs or Tyranid swarms. Thus, even if it lacks the reach of the other variants, the Leman Russ Exterminator can decimate infantry in the blink of an eye, before they reach the Astra Militarum's lines.

Its Twin Autocannons have the \textbf{Turret} ability, and the Leman Russ Exterminator has the \textbf{Slave (Hunt)} ability.

\begin{unitstatstable}
\SetCell[r=2]{c,m} 15 cm &
\SetCell[r=2]{c,m} 2+ &
\SetCell[r=2]{c,m} +3 &
\SetCell[r=2]{c,m} 8 &
Twin Autocannons &
45 cm &
2 &
4+ &
$-1$
\\
& & & &
Heavy Bolters &
45 cm &
2 &
5+ &
$-1$
\end{unitstatstable}

\unitimage{leman-russ-exterminator.png}

\unitdivider

\subsection{Limousine}

The leaders of Genestealer Cults travel aboard spacious, armoured vehicles. On the battlefield it is not uncommon for them to board their precious limousines in order to race towards the enemy as quickly as possible. Although unarmed, these vehicles reflect the dominant place the leaders occupy within the Cults.

The Limousine has the \textbf{Transport (3)}, \textbf{Attached Transport}, and \textbf{Slave (Devastation)} abilities.

\begin{unitstatstable}
25 cm & 5+ & +0 & Attached & -- & -- & -- & -- & --
\end{unitstatstable}

\unitimage{limousine.png}

\unitdivider

\subsection{Medusa}

Unlike Imperial artillery pieces that project their shells over enemy walls, the Medusa sends its own directly against the ramparts in question, turning them into rubble and opening a breach so that infantry can mount the assault. Long sieges are generally accompanied by the dull report of the Medusas' siege cannons, firing day and night from entrenched and well-protected positions, and once they have carved a path through the walls and the area has been secured, the Medusas will advance to support the assault, flattening several buildings with each shell until, street by street, the city is nothing more than a field of ruins and the enemy has nowhere left to hide.

It has the \textbf{Artillery}, \textbf{Slave (Nest)}, and \textbf{Damages Buildings (AP $-4$/1)} abilities.

\begin{unitstatstable}
15 cm & 4+ & +0 & 8 & Siege Cannon & 150 cm & 1 & 3+ & $-1$
\end{unitstatstable}

\unitimage{medusa.png}

\unitdivider

\subsection{Leman Russ Punisher}

The Punisher exchanges its anti-tank capabilities for the ability to massacre infantry. It is based on the standard Leman Russ pattern, but with a number of modifications that allow it to fulfil a vital heavy anti-infantry role which, according to some, has long been lacking among Imperial Guard regiments. Like the Leman Russ Executioner and the Demolisher, it is more heavily armoured to protect itself against point-blank attacks, but where the Demolisher was created to destroy enemy fortifications at range, the Leman Russ Punisher is more effective at destroying enemy infantry at close quarters, especially during attacks orchestrated by massive armies such as Ork and Tyranid xenos.

Its Punisher Cannon has the \textbf{Turret} ability, while its Flamer has the \textbf{Reduces Cover ($-3$)} ability, and the Punisher has the \textbf{Slave (Hunt)} ability.

\begin{unitstatstable}
\SetCell[r=2]{c,m} 15 cm &
\SetCell[r=2]{c,m} 2+ &
\SetCell[r=2]{c,m} +3 &
\SetCell[r=2]{c,m} 8 &
Punisher Cannon &
20 cm &
3 &
4+ &
0
\\
& & & &
Flamer &
20 cm &
2 &
4+ &
0
\end{unitstatstable}

\unitimage{leman-russ-punisher.png}

\unitdivider

\subsection{Thunderbolt}

The Thunderbolt is the single-seat heavy fighter on the Imperium's front lines. It is the mainstay of the Aeronautica Imperialis's wings, a versatile, robust, and well-armed fighter with a good top speed and good manoeuvrability. It is greatly appreciated by its crews and has given good service for many centuries. Durable and reliable in combat, the Thunderbolt's principal role is that of an air-superiority fighter, required to seek out and engage enemy aircraft in violent combats or to hunt opposing bombers in order to establish air superiority over the battlefield. One of the Thunderbolt's strengths is its versatility, with good performance as a high-altitude escort fighter or low-altitude fighter-bomber, as a night fighter, or even as a reconnaissance aircraft.

The Thunderbolt fighter has the \textbf{Integral Armour}, \textbf{Dodge (5+)}, \textbf{Flyer}, and \textbf{Interceptor} abilities.

\begin{unitstatstable}
\SetCell[r=3]{c,m} -- &
\SetCell[r=3]{c,m} 3+ &
\SetCell[r=3]{c,m} +3 &
\SetCell[r=3]{c,m} 7 &
Twin Autocannon &
30 cm &
2 &
3+ &
0
\\
& & & &
Twin Lascannon &
45 cm &
1 &
4+ &
$-1$
\\
& & & &
Hellstrike Missile &
45 cm &
2 &
5+ &
$-2$
\end{unitstatstable}

\unitimage{thunderbolt.png}

\unitdivider

\subsection{Vulture}

The Vulture is closely related to the Valkyrie. Whereas the Valkyrie is a transport aircraft intended to deliver troops and equipment into the combat zone, the Vulture is a multi-role gunship, replacing the transport compartment in favour of a large weapons load. Like the Valkyrie, it is much slower than a Lightning or a Thunderbolt, but it carries a far greater weapons load for its size. Its ability to hover over the battlefield means that it is always close to the action and in direct support of the ground battle. Vultures generally fly in support of Valkyrie operations, providing significant firepower when airborne troops enter play.

The Vulture has the \textbf{Integral Armour} and \textbf{Flyer} abilities.

\begin{unitstatstable}
\SetCell[r=2]{c,m} 30 cm &
\SetCell[r=2]{c,m} 4+ &
\SetCell[r=2]{c,m} +1 &
\SetCell[r=2]{c,m} 7 &
Heavy Bolter &
30 cm &
1 &
5+ &
$-1$
\\
& & & &
Autocannon &
30 cm &
2 &
4+ &
0
\end{unitstatstable}

\unitimage{vulture.png}

\unitdivider

\subsection{Termite}

A specialised transport vehicle, the Termite is capable of deploying a complete platoon of warriors onto the battlefield while bypassing enemy fortifications and sentries. It is able to dig tunnels rapidly through the densest materials at speeds comparable to those of surface transport engines, thanks to its mill, which uses a set of flux-cutters and phase-shield generators to assist its progress. Once at its target location, it emerges at the surface, or even directly into the lower level of macro-fortifications, scattering nearby enemies with its fusion cutters before disgorging its deadly cargo.

It has the \textbf{Deep Strike (2)}, \textbf{Attached Transport}, \textbf{Transport (2)}, and \textbf{Slave (Hunt)} abilities.

\begin{unitstatstable}
15 cm & 4+ & +1 & Attached & Autocannon & 45 cm & 1 & 4+ & 0
\end{unitstatstable}

\unitimage{termite.png}
